% Part: computability
% Chapter: computability-theory
% Section: rice-theorem

\documentclass[../../../include/open-logic-section]{subfiles}

\begin{document}

\olfileid[tr]{cmp}{thy}{rce}
\olsection{Rice Teoremi}

Düşündüğünüzde $\fn{Tot}$'un özel ayrıntılarının
\olref[tot]{prop:total} ispatında rol oynamadığını görürsünüz. $h(x,y)$
fonksiyonunu tam olarak $x\in K$ olduğunda sabit $j(y)=0$ fonksiyonu gibi
davranacak biçimde tasarladık; fakat o koşullarda herhangi başka bir kısmi
hesaplanabilir fonksiyon gibi davranmasını da sağlayabilirdik. Bu gözlem,
kabaca hesaplanabilir fonksiyonların önemsiz olmayan hiçbir özelliğinin karar
verilebilir olmadığını söyleyen daha genel bir teorem kurmamızı sağlar.

$\cfind{0},\cfind{1},\cfind{2},\dots$ dizisinin kısmi hesaplanabilir
fonksiyonlar için standart numaralandırmamız olduğunu unutmayın.

\begin{thm}[Rice Teoremi]
  $C$, kısmi hesaplanabilir fonksiyonların herhangi bir kümesi ve
  $A=\Setabs{n}{\cfind{n}\in C}$ olsun. $A$ hesaplanabilirse ya $C$ boştur
  ya da $C$, bütün kısmi hesaplanabilir fonksiyonların kümesidir.
\end{thm}

Bir \emph{indis kümesi}, aynı fonksiyonu ``hesaplayan'' $n$ ile $m$
indislerinin ya ikisini birden içeren ya da ikisini de içermeyen bir $A$
kümesidir. Teoremdeki $A$'nın bu özelliği taşıdığını görmek zor değildir.
Tersine $A$ bir indis kümesi ve $C$ bu indislerin hesapladığı fonksiyonlar
kümesiyse $A=\Setabs{n}{\cfind{n}\in C}$ olur.

\begin{explain}
Bu terminolojiyle Rice teoremi, önemsiz olmayan hiçbir indis kümesinin karar
verilebilir olmadığını söyler. Teoremi anlamak için \emph{programlar}, örneğin
sevdiğiniz programlama dilindeki programlar ile bunların hesapladığı
fonksiyonlar arasındaki ayrımı vurgulamak yararlıdır. Sözdizimsel nesneler
olan programlar, yani indisler hakkında hesaplanabilir sorular elbette
vardır: Bu programda $150$'den fazla simge var mı? $22$'den fazla satır var
mı? Bir ``while'' deyimi var mı? ``hello world'' dizisi bir ``print''
deyiminin argümanı olarak geçiyor mu? Rice teoremi, programın
\emph{davranışı} hakkındaki önemsiz olmayan hiçbir sorunun hesaplanabilir
olmadığını söyler. Şunlar bu tür sorulardır: Program $0$ girdisinde duruyor
mu? Hiç duruyor mu? Hiç çift sayı çıktı olarak veriyor mu?
\end{explain}

\begin{proof}[Rice teoreminin ispatı]
$C$ ne boş ne de bütün kısmi hesaplanabilir fonksiyonların kümesi olsun ve
$A$, $C$ içindeki fonksiyonların indisleri kümesi olsun. $A$ hesaplanabilir
olsaydı durma problemini çözebileceğimizi göstereceğiz; böylece $A$
hesaplanabilir değildir.

Genelliği kaybetmeden hiçbir yerde tanımlı olmayan $f$ fonksiyonunun $C$
içinde olmadığını varsayabiliriz; aksi durumda aşağıdaki savda $C$ ile
tümleyeninin yerlerini değiştiririz. $g$, $C$ içindeki herhangi bir fonksiyon
olsun. Düşünce şudur: $A$'ya karar verebilseydik $f$'yi hesaplayan indislerle
$g$'yi hesaplayan indisleri ayırt edebilir ve bu yetiyi durma problemini
çözmek için kullanabilirdik.

Şöyle yaparız. Evrensel kısmi hesaplanabilir fonksiyonları kullanarak
\[
  h(x,y)\simeq\begin{cases}
    \text{tanımsız} & \text{$\cfind{x}(x)\fundefined$ ise}\\
    g(y) & \text{aksi durumda}
  \end{cases}
\]
fonksiyonunu tanımlayabiliriz. $h$'yi hesaplamak için önce
$\cfind{x}(x)$ değerini hesaplamayı deneriz. Bu hesap durursa $g(y)$
hesabına geçer, \emph{o} hesap da durursa çıktıyı döndürürüz. Daha biçimsel
olarak
\[
  h(x,y)\simeq\Proj{2}{0}(g(y),\fn{Un}(x,x))
\]
yazabiliriz; burada $\Proj{2}{0}(z_0,z_1)=z_0$, hesaplanabilir olan ve
$0$. argümanı döndüren iki yerli izdüşüm fonksiyonudur.

$h$ kısmi hesaplanabilir fonksiyonların bir bileşkesidir ve sağ taraf tam
olarak $\fn{Un}(x,x)$ ile $g(y)$ değerlerinin ikisi de tanımlı olduğunda
tanımlı ve $g(y)$'ye eşittir.

Sabit bir $x$ için $\cfind{x}(x)$ tanımsızsa $h(x,y)$ her $y$ için
tanımsızdır; $\cfind{x}(x)$ tanımlıysa $h(x,y)\simeq g(y)$ olur. Böylece
her sabit $x$ değerinde $h(x,y)$ ya $f$ ya da $g$ gibi davranır ve
$\cfind{x}(x)$'in tanımlı olup olmadığına karar vermek bu iki durumdan
hangisinin geçerli olduğuna karar vermektir. Fakat bu, $h_x(y)\simeq h(x,y)$
fonksiyonunun $C$ içinde olup olmadığına karar vermektir; $A$
hesaplanabilir olsaydı bunu yapabilirdik.

Daha biçimsel olarak $h$ kısmi hesaplanabilir olduğundan bazı $e$ indisi için
$\cfind{e}$ fonksiyonuna eşittir. $s$-$m$-$n$ teoremi uyarınca her $x$ için
$\cfind{s(e,x)}(y)=h_x(y)$ olacak ilkel özyinelemeli bir $s$ fonksiyonu
vardır. Artık her $x$ için $\cfind{x}(x)\fdefined$ ise
$\cfind{s(e,x)}$, $g$ ile aynı fonksiyondur ve $s(e,x)\in A$ olur. Öte
yandan $\cfind{x}(x)\uparrow$ ise $\cfind{s(e,x)}$, $f$ ile aynı
fonksiyondur ve $s(e,x)\notin A$ olur. Başka bir deyişle her $x$ için
$x\in K$ ancak ve ancak $s(e,x)\in A$ olur. $A$ hesaplanabilir olsaydı
$K$ de hesaplanabilir olurdu; bu bir çelişkidir. Dolayısıyla $A$
hesaplanabilir değildir.
\end{proof}

Rice teoremi çok güçlüdür. Aşağıdaki doğrudan sonuç bazı örnek uygulamalar
verir.
\begin{cor}
Aşağıdaki kümeler karar verilemezdir:
\begin{enumerate}
\item $\Setabs{x}{\text{$17$, $\cfind{x}$ fonksiyonunun görüntü kümesindedir}}$
\item $\Setabs{x}{\text{$\cfind{x}$ sabittir}}$
\item $\Setabs{x}{\text{$\cfind{x}$ tümeldir}}$
\item $\Setabs{x}{\text{$y<y'$ ve $\cfind{x}(y)\fdefined$ olduğunda,
    $\cfind{x}(y')\fdefined$ ise $\cfind{x}(y)<\cfind{x}(y')$ olur}}$.
\end{enumerate}
\end{cor}

\begin{proof}
  Bunların tümü önemsiz olmayan indis kümeleridir.
\end{proof}

\end{document}

